\definecolor{myred}{RGB}{200, 30, 30}

\begin{figure}[h!]
\centering
\begin{tikzpicture}[scale=0.9]

  % --- Quarks of R-baryon (top group) ---
  \draw[gluon, myred]   (-3, 2.0) -- (3, 2.0);
  \draw[myline, myred]  (-3, 2.0) -- (3, 2.0);
  \node[left, myred] at (-3, 2.0) {\(\tilde{g}\)};
  \node[right, myred] at (3, 2.0) {\(\tilde{g}\)};

  \draw[fermion] (-3, 1.5) -- (3, 1.5);
  \node[left] at (-3, 1.5) {\(d\)};
  \node[right] at (3, 1.5) {\(d\)};

  \draw[fermion] (-3, 1.0) -- (3, 1.0);
  \node[left] at (-3, 1.0) {\(d\)};
  \node[right] at (3, 1.0) {\(d\)};

  % --- u quark from R-baryon -> goes to π+ ---
  \coordinate (Vu_plus) at (1.5, 0.05);
  \coordinate (Vu_minus) at (2, -0.03);
  \draw[fermion] (-3, 0.5) -- (Vu_plus);
  \draw[myline] (Vu_plus) to[bend right=90] (Vu_minus);
  \draw[myline] (Vu_minus) -- (3, -0.1);
  \node[left] at (-3, 0.5) {\(u\)};
  \node[right] at (3, -0.1) {\(u\)};

  % --- Gluon emitted by u, splits into d dbar ---
  \coordinate (Vg) at (0.6, -0.4);
  \draw[gluon]   (-0.3,-1.7) -- (Vg);
    \node[above]   at (-0.3, -1.3) {\(g\)};

  \draw[fermion] (Vg) -- (3, 0.5);   % d -> joins R-baryon
  \node[right] at (3, 0.5) {\(d\)};
  \draw[fermion] (3,-0.6) -- (Vg);   % dbar -> joins π+
  \node[right] at (3,-0.6) {\(\overline{d}\)};

  % --- Neutron quarks (all spectators) ---
  \draw[fermion] (-3,-1.7) -- (-0.3,-1.7);
  \draw[fermion] (-0.3,-1.7) -- (3,-1.7);
  \node[left] at (-3,-1.7) {\(u\)};
  \node[right] at (3,-1.7) {\(u\)};

  \draw[fermion] (-3,-2.1) -- (3,-2.1);
  \node[left] at (-3,-2.1) {\(d\)};
  \node[right] at (3,-2.1) {\(d\)};

  \draw[fermion] (-3,-2.5) -- (3,-2.5);
  \node[left] at (-3,-2.5) {\(d\)};
  \node[right] at (3,-2.5) {\(d\)};

  % --- Initial-state braces ---
  \draw[thick, decorate, decoration={brace, amplitude=4pt, mirror}]
        (-3.5, 2.1) -- (-3.5, 0.4) node[midway, left=6pt] {\(R^{0}_{\tilde{g}udd}\)};
  \draw[thick, decorate, decoration={brace, amplitude=4pt, mirror}]
        (-3.5,-1.6) -- (-3.5,-2.6) node[midway, left=6pt] {\(n\)};
  
  % --- Final-state braces ---
  \draw[thick, decorate, decoration={brace, amplitude=4pt}]
        ( 3.5, 2.1) -- ( 3.5, 0.4) node[midway, right=6pt] {\(R^{-}_{\tilde{g}ddd}\)};
  \draw[thick, decorate, decoration={brace, amplitude=4pt}]
        ( 3.5, 0.0) -- ( 3.5,-0.7) node[midway, right=6pt] {\(\pi^{+}\)};
  \draw[thick, decorate, decoration={brace, amplitude=4pt}]
        ( 3.5,-1.6) -- ( 3.5,-2.6) node[midway, right=6pt] {\(n\)};

\end{tikzpicture}
\caption{Feynman-like diagram of an R-hadron undergoing a charge-flipping interaction with a neutron: $R^{0}_{\tilde{g}udd} + n \to R^{-}_{\tilde{g}ddd} + n + \pi^{+}$.}
\label{fig:rhadron_charge_flip}
\end{figure}